\section{Pioneer RKNN: нейросети на борту}
\label{sec:rknn}

\begin{infobox}[Цели главы]
\begin{itemize}
    \item Понимать, что такое нейросеть-детектор и зачем дрону нейропроцессор.
    \item Находить объекты, позу человека и текст на кадрах камеры библиотекой \texttt{pioneer\_rknn}.
    \item Строить конвейер «камера $\to$ нейросеть $\to$ решение $\to$ команда дрону».
    \item Собрать автономную миссию: дрон сам ищет человека и сообщает о находке.
\end{itemize}
\end{infobox}

\subsection{Нейросеть и нейропроцессор}

В главе~\ref{sec:python-vision} дрон находил на кадре маркеры: у маркера строгая форма, и найти его можно правилами. Человека, машину или дорожный знак правилами не описать --- их находит \textbf{нейросеть}. Для программы это просто функция: на входе кадр, на выходе список найденных объектов с рамками и уверенностью. Как распознавать, нейросеть «выучила» заранее на тысячах размеченных картинок; результат обучения хранится в файле модели.

Нейросети требуют много вычислений, и обычный процессор дрона справляется с ними медленно. В процессоре Rockchip у Пионера Мини~2 и Radxa Zero~3W есть отдельный блок --- \textbf{нейропроцессор} (NPU), который считает нейросети в разы быстрее. Модель для него переводится в особый формат --- файл \texttt{.rknn}.

Библиотека \texttt{pioneer\_rknn} прячет работу с NPU за несколькими классами. Как и \texttt{pioneer\_sdk2}, она встроена в образ Pioneer~OS и работает только на борту; в симуляторе её нет. Программы этой главы используют обе библиотеки: \texttt{pioneer\_sdk2} даёт кадры и полёт, \texttt{pioneer\_rknn} --- распознавание.

\begin{center}
\small
\begin{tabularx}{\linewidth}{@{}>{\ttfamily}l X >{\raggedright\arraybackslash}p{0.27\linewidth}@{}}
\toprule
\textrm{\textbf{Класс}} & \textbf{Что делает} & \textbf{Архитектуры} \\
\midrule
Yolo & находит объекты: рамка, класс, уверенность & \texttt{yolov8}, \texttt{yolov11} \\
YoloPose & находит людей и 17 точек скелета & \texttt{yolov8-pose} \\
PaddleOCR & находит строки текста и читает их & PP-OCRv5 \\
ModelContainer & загрузка модели и запуск на NPU; основа остальных классов & любая \texttt{.rknn} \\
ModelRegistry & список моделей на дроне, загрузка новых & --- \\
\bottomrule
\end{tabularx}
\end{center}

\subsection{Какие модели есть на дроне}

Модели хранятся в \textbf{регистре} --- службе на борту дрона. По имени из регистра модель и загружают: \verb|Yolo(model_name="yolov8n")|. Другой способ --- путь к своему файлу: \verb|Yolo(model_path="/home/pioneermini/best.rknn")|.

\begin{codefile}[python]{python-onboard/pioneer_rknn/rknn_01_models.py}
\end{codefile}

\paragraph{Что происходит.} \verb|list_model()| возвращает словарь с моделями регистра, \verb|get_model_info(имя)| --- сведения об одной модели; при ошибке обе возвращают \verb|None|. Запустите эту программу первой: имена моделей в примерах ниже должны совпадать с теми, что она напечатает.

\subsection{Поиск объектов: Yolo}

YOLO (\emph{You Only Look Once}) --- семейство нейросетей, которые за один проход по кадру находят все объекты. Модель \texttt{yolov8n} обучена на наборе COCO и различает 80 классов: человек (класс 0), машина (2), собака (16), бутылка (39) и другие. Буква \texttt{n} (\emph{nano}) --- самый маленький и быстрый вариант.

\begin{codefile}[python]{python-onboard/pioneer_rknn/rknn_02_detect.py}
\end{codefile}

\paragraph{Что происходит.} Цикл --- это конвейер из четырёх шагов:
\begin{enumerate}
    \item \textbf{Кадр.} \verb|camera.get_cv_frame()| --- изображение BGR, например $640\times480$.
    \item \textbf{Подготовка.} Нейросеть принимает вход строго заданной формы. \verb|prepare| сжимает кадр до $640\times640$, переставляет цвета в RGB (модели обучали на RGB) и добавляет измерение «пакет»: \verb|np.expand_dims| превращает массив формы \verb|(640, 640, 3)| в \verb|(1, 640, 640, 3)| --- пакет из одной картинки.
    \item \textbf{Нейросеть.} \verb|model.run([...])| считает модель на NPU и отбирает найденное (\emph{постобработка}).
    \item \textbf{Результат.} Рамки пересчитываются из координат входа $640\times640$ в координаты кадра (\verb|kx|, \verb|ky|) и рисуются; кадр уходит в браузер, как в главе~\ref{sec:sdk2}.
\end{enumerate}

\begin{hintbox}[Формат ответа --- в одной функции]
\verb|Yolo.run| возвращает кортеж (рамки, классы, уверенности): рамка --- \verb|(x1, y1, x2, y2)| в пикселях входа $640\times640$; если объектов нет --- \verb|None|. Так устроена постобработка YOLOv8 в примерах Rockchip (\texttt{rknn\_model\_zoo}). Весь разбор ответа собран в функции \verb|detections|: если в вашей версии образа ответ выглядит иначе, напечатайте его (\verb|print(result)|) и поправьте только эту функцию --- остальная программа от формата не зависит.
\end{hintbox}

\subsubsection*{Пороги}

Нейросеть не говорит «да» или «нет» --- она оценивает уверенность от 0 до 1. Аргумент \verb|object_thresh| конструктора (по умолчанию 0.25) отбрасывает всё, в чём модель уверена меньше. Один и тот же объект модель часто находит несколькими перекрывающимися рамками; \textbf{NMS} (\emph{non-maximum suppression}) оставляет из них самую уверенную, а \verb|nms_thresh| (0.45) задаёт, насколько сильно рамки должны перекрываться, чтобы считаться одним объектом. Повысили \verb|object_thresh| --- меньше ложных находок, но и настоящие объекты теряются чаще.

\begin{rulebox}[Освобождайте нейропроцессор]
\verb|model.release()| возвращает ресурсы NPU. Вызывайте её в \verb|finally|, как \verb|close_connection()|: тогда модель освобождается и после \key{Ctrl}+\key{C}, и после ошибки.
\end{rulebox}

\subsection{Поза человека: YoloPose}

\verb|YoloPose| находит людей и для каждого --- 17 ключевых точек скелета. В главе~\ref{sec:human-tracking} это делала библиотека MediaPipe на компьютере; здесь то же самое считается на борту.

\verb|run| возвращает список найденных людей. У каждого есть \verb|score| (уверенность), рамка \verb|xmin|, \verb|ymin|, \verb|xmax|, \verb|ymax| и \verb|keypoint| --- массив $17\times3$: для каждой точки $x$, $y$ и уверенность. Номера точек --- по стандарту COCO:

\begin{center}
\small
\begin{tabular}{@{}rl rl rl@{}}
\toprule
0 & нос & 5, 6 & плечи (лев., прав.) & 11, 12 & бёдра \\
1, 2 & глаза & 7, 8 & локти & 13, 14 & колени \\
3, 4 & уши & 9, 10 & запястья & 15, 16 & лодыжки \\
\bottomrule
\end{tabular}
\end{center}

\begin{codefile}[python]{python-onboard/pioneer_rknn/rknn_03_pose_signal.py}
\end{codefile}

\paragraph{Что происходит.} Жест «руки вверх» --- обе кисти выше плеч. Ось $y$ на картинке направлена \emph{вниз}, поэтому «выше» значит «меньше по $y$». Подсветка меняется только при смене жеста (\verb|now != signal|): незачем двадцать раз в секунду отправлять автопилоту одну и ту же команду. Дрон при этом не взлетает --- жест только включает зелёный или красный цвет, так программу безопасно проверить, держа дрон в руках.

\subsection{Чтение текста: PaddleOCR}

Распознавание текста (OCR) --- две нейросети подряд. Первая (\emph{детектор}) находит на кадре строки текста и обводит их четырёхугольниками, вторая (\emph{распознаватель}) читает буквы внутри каждого.

\begin{codefile}[python]{python-onboard/pioneer_rknn/rknn_04_ocr.py}
\end{codefile}

\paragraph{Что происходит.} В отличие от \verb|Yolo|, \verb|ocr.run| получает кадр целиком, без подготовки, и возвращает два списка: четырёхугольники строк и варианты прочтения для каждой строки --- пары (текст, уверенность). Программа берёт первый вариант и отбрасывает неуверенные. Названия моделей должны совпадать с именами в регистре (\verb|rknn_01_models.py|).

\subsection{Проект: найти человека}

Соберём всё вместе. Дрон взлетает, медленно поворачивается на месте и смотрит в камеру; увидев человека, останавливается, включает зелёный свет и садится. Не нашёл за 30~секунд --- красный свет и посадка.

\begin{codefile}[python]{python-onboard/pioneer_rknn/rknn_05_search_person.py}
\end{codefile}

\paragraph{Что происходит.} Это конечный автомат из трёх состояний: \verb|SEARCH| $\to$ \verb|FOUND| или \verb|NOT_FOUND|, затем посадка. Три приёма делают его надёжным:
\begin{itemize}
    \item \textbf{Подтверждение.} Нейросеть иногда ошибается на одном кадре. Человек засчитывается, только если его видно \verb|CONFIRM| кадров подряд: счётчик \verb|seen| растёт, пока человек виден, и обнуляется, как только пропал.
    \item \textbf{Короткий interval.} Поворот задаётся на каждом кадре с \verb|interval=0.5|. Если программа зависнет, через полсекунды дрон перестанет вращаться сам.
    \item \textbf{Таймаут.} Поиск не бесконечен: через \verb|SEARCH_TIME| секунд автомат уходит в \verb|NOT_FOUND|.
\end{itemize}

\begin{warnbox}[Первый запуск --- без людей рядом]
Проверяйте программу сначала на земле (закомментируйте взлёт) и с человеком на фото, а не в двух метрах от винтов. Дрон, который «ищет человека», не должен лететь к нему.
\end{warnbox}

\subsection{Своя модель}

Модель из регистра знает 80 классов COCO. Чтобы дрон находил что-то своё --- дорожные знаки, посадочные метки, детали на конвейере, --- модель дообучают:
\begin{enumerate}
    \item собирают и размечают фотографии нужных объектов;
    \item обучают YOLOv8 или YOLO11 на компьютере или в облаке (библиотека Ultralytics);
    \item переводят модель в формат \texttt{.rknn} утилитой RKNN-Toolkit2 под процессор дрона;
    \item копируют файл на дрон и загружают: \verb|Yolo(model_path="best.rknn")| --- или добавляют в регистр командой \verb|upload_model| класса \verb|ModelRegistry|: имя, версия, путь к файлу и архитектура, например \verb|"yolov8"|.
\end{enumerate}
У своей модели свои классы: номера в ответе \verb|run| --- это номера из вашего набора данных, а не COCO. Готовые модели для NPU Rockchip собраны в репозитории \url{https://github.com/airockchip/rknn_model_zoo}.

\subsection{Задания к главе}

\subsubsection*{Уровень 1}
\begin{enumerate}
    \item Измерьте скорость конвейера: сколько кадров в секунду обрабатывает \verb|rknn_02_detect.py|? Выводите число на кадр.
    \item Измените детектор так, чтобы он показывал только людей и бутылки.
\end{enumerate}

\subsubsection*{Уровень 2}
\begin{enumerate}
    \item Посчитайте людей в кадре и показывайте число цветом подсветки: никого --- выключена, один --- зелёный, больше --- синий.
    \item Добавьте в \verb|rknn_03_pose_signal.py| второй жест: поднята только одна рука.
\end{enumerate}

\subsubsection*{Уровень 3: мини-проект}
\textbf{Следование за человеком.} Дрон поворачивается так, чтобы человек был в центре кадра: \verb|yaw_rate| пропорционален смещению центра рамки от центра кадра (P-регулятор из главы~\ref{sec:human-tracking}). Не забудьте про короткий \verb|interval|, ограничение $|yaw\_rate| \le 1.82$ и посадку, если человек потерян дольше 5~секунд.

\subsection*{Итоги}
\begin{itemize}
    \item Нейросеть --- функция «кадр $\to$ объекты»; NPU считает её быстро, модель для него --- файл \texttt{.rknn}.
    \item \verb|Yolo| --- объекты, \verb|YoloPose| --- скелет из 17 точек, \verb|PaddleOCR| --- текст; модели берутся из регистра по имени или из файла.
    \item Вход Yolo готовят: $640\times640$, RGB, пакет из одного кадра; рамки пересчитывают обратно в координаты кадра.
    \item Разбор ответа модели держите в одной функции; результат одного кадра подтверждайте несколькими.
    \item \verb|model.release()| --- в \verb|finally|.
\end{itemize}
